% ---------------------------
% CONFIGURATION
% ---------------------------
\documentclass[11pt,a4paper]{moderncv}
\moderncvstyle{classic}
\moderncvcolor{black}
\definecolor{firstnamecolor}{rgb}{0,0,0}

\usepackage[utf8]{inputenc}
\usepackage[T1]{fontenc}
\usepackage[scale=0.9]{geometry}
\usepackage[dvipsnames]{xcolor}
\definecolor{PaleBlue}{RGB}{76, 151, 215}
\usepackage[colorlinks=true, urlcolor=PaleBlue, linkcolor=PaleBlue]{hyperref}
\usepackage{microtype}

\renewcommand{\rmdefault}{ppl}
\renewcommand{\sfdefault}{phv}
\renewcommand*{\namefont}{\fontsize{26}{18}\mdseries}

\microtypesetup{expansion=false}

\input{../../secrets/personal_info_dk.tex}

\title{Curriculum Vitae}

% ---------------------------

\begin{document}
\makecvtitle

% ---------------------------
% Introduction
% ---------------------------
\section*{Introduction}

Research-oriented ML engineer with seven years of experience in computer vision, representation learning, probabilistic modelling, and production ML.
\newline\hspace*{2em}
My research includes learned representations of segmented 3D CT scans and spatial dose prediction. Independently, I have built object-detection and visual-embedding systems for real-world imagery.
\newline\hspace*{2em}
Strong mathematical background with hands-on experience designing, implementing, debugging, and evaluating deep-learning systems in Python.

% ---------------------------
% Education
% ---------------------------
\section{Education}

\cventry{2015--2020}{Royal Institute of Technology, KTH -- Data Science and Engineering}{}{GPA: 5.0 out of 5.0}{}{Master's specialisation in Statistical Learning and Data Analytics; Bachelor's degree in Engineering Physics.
\href{https://kth.diva-portal.org/smash/get/diva2:1431327/FULLTEXT02.pdf}{Thesis: Image-Distance Learning for Probabilistic Spatial Dose Prediction in Radiation Therapy}.
\newline
Relevant coursework: Statistical Machine Learning; Modern Methods of Statistical Learning; Reinforcement Learning; Optimisation; Control Theory.}

\cventry{2019}{Swiss Federal Institute of Technology, ETH}{Exchange in Applied Mathematics}{}{}{}

\cventry{2017--2022}{Stockholm University}{Supplementary Bachelor's Degree in Economics}{}{}{}

% ---------------------------
% Selected Experience
% ---------------------------
\section{Selected Experience}

% ---------------------------
\large{Researcher within Data Science @ RaySearch Laboratories}
\normalsize
\begin{itemize}
    \item Developed learning-based methods for automated radiotherapy planning from segmented 3D CT scans.
    \item Learned latent anatomical representations with autoencoders for spatial dose prediction.
    \item Designed and evaluated methods around clinical constraints with medical physicists and clinicians.
\end{itemize}

{\footnotesize\textit{Keywords: 3D medical imaging, representation learning, deep learning, spatial prediction, probabilistic modelling}}

\hfill{\small{\textit{\hyperref[sec:raysearch]{Read more...}}}}

% ---------------------------
\large{Principal Data Scientist @ Valcon}
\normalsize
\begin{itemize}
    \item Led ML projects from problem formulation and experimental evaluation through deployment.
    \item Built a real-time deep neural network with a custom training pipeline and modular architecture.
    \item Owned implementation, debugging, testing, performance analysis, and production integration.
\end{itemize}

{\footnotesize\textit{Keywords: deep learning, experimental evaluation, Python, production ML, software engineering}}

\hfill{\small{\textit{\hyperref[sec:valcon]{Read more...}}}}

% ---------------------------
\large{Solo Project @ Project \texttt{iris}}
\normalsize
\begin{itemize}
    \item Built an end-to-end vision system for object detection and visual retrieval in unconstrained images.
    \item Implemented OpenCLIP embeddings, vector search, backend APIs, web scraping, and a Chromium extension.
\end{itemize}

{\footnotesize\textit{Keywords: computer vision, object detection, OpenCLIP, representation learning, vector search}}

\hfill{\small{\textit{\hyperref[sec:iris]{Read more...}}}}

% ---------------------------
% Skills
% ---------------------------
\section{Skills and Other}

\cvitem
{Machine Learning}
{Deep learning, computer vision, representation learning, object detection, 3D medical imaging, probabilistic modelling, model evaluation.}

\cvitem
{Technology}
{Python, PyTorch, TensorFlow, OpenCLIP, YOLOv8, scikit-learn, Linux, Git, Docker, FastAPI.}

\cvitem
{Soft-Skills}
{Independent problem solving, communication, interdisciplinary collaboration, mentoring.}

\newpage

% ---------------------------
% Detailed Work Experience
% ---------------------------
\section{Detailed Work Experience}

% ---------------------------
\phantomsection\label{sec:signum}
\cventry{2024--Present}{Data Scientist / ML Engineer}{Signum Life Science}{Copenhagen}{}{
At Signum, I work across machine learning, software engineering, architecture, and project ownership, turning loosely defined pharmaceutical analytics problems into maintainable production systems.
\newline\hspace*{2em}
I help design and build an end-to-end forecasting platform covering data validation, feature engineering, model training, experiment tracking, registration, and batch inference. The architecture supports thousands of independently trained models.
\newline\hspace*{2em}
My work includes failure analysis, reproducible evaluation, and ensuring experimental improvements generalise to heterogeneous real-world data. Python, SQL, Databricks, and MLflow form the core implementation.
}{}

% ---------------------------
\phantomsection\label{sec:valcon}
\cventry{2022--2024}{Principal Data Scientist}{Valcon}{Copenhagen}{}{
At Valcon, I led machine-learning projects across manufacturing, pharma, energy, IoT, and life science, covering technical scoping, modelling, experimental evaluation, software engineering, and deployment.
\newline\hspace*{2em}
For one project, I designed and deployed a real-time deep neural network with natural-language inputs and high-dimensional targets. I built the training pipeline and modular model architecture and owned debugging, testing, inference, and performance optimisation.
\newline\hspace*{2em}
I regularly worked with noisy and sensor-derived data, where careful failure analysis and iterative experimentation were essential. I also mentored colleagues and delivered training in explainable AI and Git.
}{}

% ---------------------------
\phantomsection\label{sec:valtech}
\cventry{2021--2022}{Data Scientist}{Valtech}{Copenhagen}{}{
Delivered applied ML solutions, including forecasting, customer segmentation, and behavioural analysis. Worked closely with business stakeholders to translate outputs into commercial decisions.
}{}

% ---------------------------
\phantomsection\label{sec:raysearch}
\cventry{2019--2021}{Researcher within Data Science}{RaySearch Laboratories}{Stockholm}{}{
At RaySearch Laboratories, I researched automated radiotherapy planning using deep learning, probabilistic modelling, and optimisation on high-dimensional 3D medical data.
\newline\hspace*{2em}
I learned latent anatomical representations from segmented 3D CT scans with autoencoders, used similarity metrics to retrieve related patients, and predicted dose-volume histograms and full spatial dose distributions with sparse Gaussian processes. The work formed my master's thesis and continued afterwards.
\newline\hspace*{2em}
I also redesigned lexicographic optimisation methods with clinicians to better represent clinical priorities.
}{}

% ---------------------------
% Projects
% ---------------------------
\section{Projects}

\cvitem{\phantomsection\label{sec:iris}Project \texttt{iris}}{
An end-to-end computer-vision retrieval system for unconstrained e-commerce imagery. I built the complete pipeline: object detection, OpenCLIP embeddings, vector indexing, similarity search, backend APIs, persistence, dynamic web scraping, and a Chromium extension.
\newline\hspace*{2em}
The project combined pretrained vision models, representation learning, retrieval, and software engineering on diverse real-world images.
}

\cvitem{\phantomsection\label{sec:resume_builder}\texttt{resume\_ builder3}}{
A source-grounded agentic workflow for tailored job applications using structured profile data, job listings, LaTeX templates, and evidence mappings.
\newline\hspace*{2em}
I built the human-in-the-loop workflow to crawl listings, create digests and GitHub issues, generate application branches, render documents, and open pull requests for review.
}

\end{document}
